\begin{problem}[81.]
Recall that given a complex number $z = a + bi$ it's complex conjugate, denoted $\conj{z}$, is given by $\conj{z} = a - bi$.

\begin{enumerate}[(a)]
    \item Prove that $|\conj{z}| = |z|$
    \item Prove that $|z| = \sqrt{z \cdot \conj{z}}$
    \item Show that $Re(z) = \frac{z + \conj{z}}{2}$ and $Im(z) = \frac{z - \conj{z}}{2i}$
    \item Show that if $\theta \in arg(z)$ then $-\theta \in arg(\conj{z})$. Interpret this result geometrically. 
\end{enumerate} 
\end{problem}

\begin{solution}[(a) Prove that $|\conj{z}| = |z|$]

    \begin{align}
        z = a + bi && \conj{z} = a - bi \\
        Re(z) = a && Re(\conj{z}) = a \\
        Im(z) = b && Im(\conj{z}) = -b
    \end{align}
    
    \[|z| = \sqrt{Re(z)^2 + Im(z)^2} = \sqrt{a^2 + b^2}\]
    \[|\conj{z}| = \sqrt{Re(\conj{z})^2 + Im(\conj{z})^2} = \sqrt{a^2 + (-b)^2}\]
    \[\sqrt{a^2 + (b)^2} = \sqrt{a^2 + (-b)^2}\]
    \[|z| = |\conj{z}|\]
\end{solution}

\hrule \medskip 

\begin{solution}[(b) Prove that $|z| = \sqrt{z \cdot \conj{z}}$]
    \begin{align}
        \sqrt{z \cdot \conj{z}} &= \sqrt{(a + bi) \cdot (a - bi)} \\
        &= \sqrt{a^2 - abi + abi - bi^2} \\
        &= \sqrt{a^2 - bi^2} \\
        i^2 &= -1 \\
        \sqrt{z \cdot \conj{z}} &= \sqrt{a^2 + b^2} = |z|
    \end{align}
\end{solution}

\hrule \medskip 

\begin{solution}[(c) Show that $Re(z) = \frac{z + \conj{z}}{2}$ and $Im(z) = \frac{z - \conj{z}}{2i}$]
    \begin{align}
        Re(z) &= \frac{z + \conj{z}}{2} \\
        &= \frac{(a+bi) + (a - bi)}{2} \\
        &= \frac{2a}{2} = a \\
        Re(z) &= a = \frac{z + \conj{z}}{2} 
    \end{align}
    
    \pagebreak 
    
    \begin{align}
        Im(z) &= \frac{z - \conj{z}}{2i} \\
        &= \frac{(a+bi) - (a - bi)}{2i} \\
        &= \frac{2bi}{2i} = b \\
        Im(z) &= b = \frac{z - \conj{z}}{2i} 
    \end{align}
\end{solution}

\hrule \medskip 

\begin{solution}[(d) Show that if $\theta \in arg(z)$ then $-\theta \in arg(\conj{z})$. Interpret this result geometrically.]
    . \linebreak 
    
    For this we'll let $\phi$ be an angle for $\conj{z}$ and $\theta$ will be an angle for $z$.
    
    \[Im(\conj{z}) = -Im(z)\]
    \[Re(\conj{z}) = Re(z)\]
    
    \[\theta = arctan \left( \frac{Im(z)}{Re(z)} \right) \]
    \[ \phi = arctan \left( \frac{Im(\conj{z})}{Re(\conj{z})} \right) \] 
    
    Due to the above relationships this can otherwise be written as:
    
    \[ \phi = arctan \left( \frac{-Im(z)}{Re(z)} \right) \]
    
    Therefore: 
    
    \[\phi = -\theta\]
    
    I'm not 100\% sure where to go from here. How do I fit this into the argument format? I know this final step is incredibly simple but it just isn't quite clicking for me. 
\end{solution}